\section{Existence and uniqueness of the solved rate}
\label{sec:existence}

This section establishes that the fixed-point equation
\eqref{eq:fixed-point} has exactly one solution on the search bracket and
that an on-chain bisection finds it deterministically. The proof rests on
two lemmas:

\begin{itemize}
  \item \textbf{Lemma~A} (\S\ref{sec:existence:lemma-a}): the conservative
    BPT valuation $P_A$ is non-decreasing and Lipschitz in the senior
    rate $y$. The argument is the envelope theorem and requires no
    pool-specific algebra.
  \item \textbf{Lemma~B} (\S\ref{sec:existence:lemma-b}): the waterfall
    $F$ is non-decreasing and 1-Lipschitz in $\JTraw$. The argument is
    structural and does not need $F$ to be differentiable.
\end{itemize}
Composing the two lemmas (\S\ref{sec:existence:composition}) gives a
strict monotonicity bound on $E$. The conditions auditors must verify per
market are collected in \S\ref{sec:existence:obligations}; their robustness
across pool types is discussed in \S\ref{sec:existence:generality}.

\subsection{The property to prove}
\label{sec:existence:property}

We want exactly one statement, strong enough to imply both root-existence
and bisection-convergence, that pure monotonicity arguments can establish.

\begin{proposition}\label{prop:E}
On the search bracket $[x_{\mathrm{lo}}, x_{\mathrm{hi}}]
= [N/\beta,\,N/\alpha]$, the error function
$E(x) = F\!\left(P_A(x/N)\right) - x$ is continuous and strictly
decreasing. If, in addition, $E(x_{\mathrm{lo}}) \geq 0$ and
$E(x_{\mathrm{hi}}) \leq 0$, then $E$ has a unique root $x^*$ in the
bracket and bisection converges to it.
\end{proposition}

The sign-change requirement is checked at runtime by the solver
(\S\ref{sec:solver}); failure to straddle the root is treated as a market
condition the pool cannot express, and the sync reverts rather than
publishing a stale rate.

\subsection{Lemma~A: $P_A$ is monotone and Lipschitz in the rate}
\label{sec:existence:lemma-a}

Fix the E-CLP invariant $L$ implied by the current pool. The conservative
pool TVL at senior rate $y$ is
\[
  V(y) \;=\; x(y) \;+\; y\,\mathrm{yb}(y),
\]
where $\mathrm{yb}(y)$ is the senior-share reserve at the point on the
ellipse whose marginal price is $y$, and $x(y)$ is the quote reserve at
that same point. Moving along the curve, $dx/d(\mathrm{yb}) = -y$ (the
slope of the curve is the price), so
\begin{align*}
  \frac{dV}{dy}
    &\;=\; \underbrace{\frac{dx}{d(\mathrm{yb})}\cdot\frac{d(\mathrm{yb})}{dy}
          \;+\; y\cdot\frac{d(\mathrm{yb})}{dy}}_{=\,0 \text{ along the curve}}
       \;+\; \mathrm{yb}(y) \\
    &\;=\; \mathrm{yb}(y) \;\geq\; 0.
\end{align*}
This is the envelope theorem: the first-order effect of the rebalancing
variable ($\mathrm{yb}$) cancels at the optimum, leaving only the direct
term. Concretely, $V$ is non-decreasing and concave in $y$, with derivative
the in-curve senior-share count.

\begin{lemma}\label{lem:PA}
$P_A$ is non-decreasing in $y$ and Lipschitz with constant
$f \cdot \sup_y \mathrm{yb}(y)$.
\end{lemma}

\begin{proof}
$P_A(y) = \lfloor\, f \cdot V(y)\,\rfloor$. The floor operator is
1-Lipschitz and non-decreasing. From $V'(y) = \mathrm{yb}(y) \geq 0$ we
get $V$ non-decreasing, hence so is $P_A$. For any $y > y'$,
\[
  |V(y) - V(y')| \;\leq\; \sup_{\eta\in[y',y]} \mathrm{yb}(\eta) \cdot (y - y')
     \;\leq\; \sup_{\eta} \mathrm{yb}(\eta) \cdot (y - y'),
\]
so $|P_A(y) - P_A(y')| \leq f \cdot \sup_y \mathrm{yb}(y) \cdot |y - y'|$,
modulo the 1-wei rounding noise of the floor.\qed
\end{proof}

The fair-reserve formula $V(y)$ for the Gyro E-CLP is derivable
symbolically from the pool's $\tau$, $\eta$, $\zeta$ and $A^{-1}$ maps;
the resulting closed form round-trips against the pool's own
\code{computePrice} to within $2 \cdot 10^{-16}$.

A useful corollary: because $V$ depends only on the invariant and the rate
--- not on the live composition --- a single swap (which conserves $L$ by
definition) leaves $V$ unchanged. This is the source of $P_A$'s
manipulation resistance: flash-loan rebalancing of the pool cannot move
the kernel's BPT valuation.

\subsection{Lemma~B: $F$ is non-decreasing and 1-Lipschitz in $\JTraw$}
\label{sec:existence:lemma-b}

Fix $\STraw$ and the last sync's checkpoint, and vary only the input
$j = \JTraw$ that enters the waterfall. The accountant computes
$F(j)$ through four branches (junior loses, junior gains, senior loses,
senior gains), each a chain of additions, subtractions, $\min$, and
$\max$.

The structural fact is that each marginal unit of $j$ reaches the senior
side through at most one channel, and through it by at most one-for-one:

\begin{itemize}
  \item if $j$ rises and the senior tranche currently carries impermanent
    loss, the marginal increase first repays \code{stImpermanentLoss}
    (slope $1$ to $\STeff$) until it is exhausted, after which any further
    increase accrues to $\JTeff$ only (slope $0$ to $\STeff$);
  \item if $j$ falls and the junior buffer is exhausted, the residual loss
    is booked against $\STeff$ (slope $1$, in the falling direction),
    until then the loss is absorbed by $\JTeff$ (slope $0$).
\end{itemize}

The other branches are mirror images. The combined effect is that
$\partial\STeff / \partial j \in \{0, 1\}$ on every branch, and the
branches join continuously at their kinks.

\begin{lemma}\label{lem:F}
$F$ is continuous, non-decreasing, and 1-Lipschitz in $\JTraw$:
\[
  0 \;\leq\; F(j) - F(j') \;\leq\; j - j' \quad \text{for } j > j'.
\]
\end{lemma}

\begin{proof}
Piecewise-affine functions with slopes in $\{0, 1\}$ joined continuously
at their kinks are continuous, non-decreasing, and have Lipschitz constant
at most their maximum slope. Differentiability is not required: the
inequality holds across the kink points by joining-continuity, so the
result extends to the whole input range.\qed
\end{proof}

This is what makes bisection the correct primitive for this problem. The
waterfall is genuinely non-smooth (the slope jumps from $0$ to $1$ at
branch boundaries); a Newton iteration would need a derivative that does
not exist at the kinks and could step across them incorrectly. Bisection
needs neither derivative nor smoothness.

\subsection{Composition: $E$ is strictly decreasing}
\label{sec:existence:composition}

Define the iteration map $T(x) = F\!\left(P_A(x/N)\right)$ from
\eqref{eq:fixed-point}. The chain $x \mapsto x/N \mapsto P_A(\cdot)
\mapsto F(\cdot)$ composes Lipschitz constants:
\[
  \frac{1}{N}
    \quad\cdot\quad
  f \cdot \sup_y \mathrm{yb}(y)
    \quad\cdot\quad
  1
    \;=\; g, \qquad
  g \;:=\; \frac{f \cdot \sup_y \mathrm{yb}(y)}{N}.
\]
$T$ is the composition of non-decreasing functions and is therefore
non-decreasing. Combining with the Lipschitz bound:
\[
  0 \;\leq\; T(x) - T(x') \;\leq\; g\,(x - x') \quad \text{for } x > x'.
\]

\begin{proposition}\label{prop:E-strict}
If $g < 1$ on the search bracket, then $E$ is strictly decreasing on it.
\end{proposition}

\begin{proof}
For $x > x'$, from the bound above,
\[
  E(x) - E(x') \;=\; \bigl[T(x) - T(x')\bigr] - (x - x')
    \;\leq\; (g - 1)\,(x - x') \;<\; 0,
\]
since $(g - 1) < 0$ and $(x - x') > 0$. Continuity of $E$ follows from
continuity of $T$ (composition of continuous functions) and the identity
map.\qed
\end{proof}

This argument never differentiates $F$, so it remains valid across the
waterfall's kinks. Together with the sign-change condition assumed in
\cref{prop:E}, $E$ has exactly one root in the bracket and bisection
converges to it.

\subsection{Per-market obligations}
\label{sec:existence:obligations}

The argument above factors into three obligations. The first two are
structural and pool-agnostic; the third is the only per-market check the
solver must perform.

\begin{description}
  \item[\textbf{O1.}] The waterfall $F$ is continuous, non-decreasing, and
    1-Lipschitz in the input $j = \JTraw$ (\cref{lem:F}). This is a
    property of the accountant code and holds for any Royco market that
    uses the standard waterfall.
  \item[\textbf{O2.}] The conservative valuation $P_A$ is non-decreasing
    and Lipschitz in $y$ with constant $f \cdot \sup_y \mathrm{yb}(y)$
    (\cref{lem:PA}). This follows from the envelope theorem and holds for
    any pool whose reserve curve is convex and downward-sloping.
  \item[\textbf{O3.}] The loop gain
    $g = f \cdot \sup_y \mathrm{yb}(y) / N$ is strictly less than $1$ on
    the chosen search bracket.
\end{description}

O3 is the only condition that depends on the live pool's geometry, the
junior tranche's BPT share, and the choice of bracket. Near the operating
rate the pool's senior-share reserve count equals its in-pool count
$b_s$, so $g \approx f \cdot b_s / N < 1$ holds robustly because the pool
cannot hold more senior shares than exist ($b_s < N$). The bound can
approach $1$ only if the bracket extends into deep-impairment rates, where
the conservative share reserve can grow sharply --- a configuration the
bracket policy (\S\ref{sec:solver}) is required to avoid.

The kernel exposes a cheap runtime check $g_{\text{proxy}} = f \cdot
\varphi$ (with $\varphi = b_s / N$) as a sanity bound on O3; the precise
relationship of $g_{\text{proxy}}$ to the true $g$ across the bracket, and
the production margin governance should configure, are discussed in
\S\ref{sec:solver}.

\subsection{Generality across pool types}
\label{sec:existence:generality}

The proof uses only that:
\begin{itemize}
  \item $F$ has slopes in $\{0, 1\}$ on each branch and is continuous
    across its kinks (structural; same accountant for any Royco market);
  \item $P_A$ is monotone and Lipschitz in the rate (envelope theorem;
    holds for any convex, downward-sloping reserve curve).
\end{itemize}

The proof does \emph{not} use a closed-form pool invariant or assume any
differentiability of $F$. The same on-chain bisection scheme therefore
applies, unchanged, to constant-product, weighted, stable, and Gyro E-CLP
pools. Only the pool-specific Lipschitz constant
$\sup_y \mathrm{yb}(y)$ changes, entering the per-market O3 check.
